\documentclass[11pt]{article}

% ---------- Packages ----------
\usepackage[a4paper,margin=1in]{geometry}
\usepackage{amsmath,amssymb,amsthm,mathtools}
\usepackage{hyperref}
\usepackage{tikz}
\usepackage{tikz-cd}
\usetikzlibrary{arrows.meta,positioning,cd}

% ---------- Theorem environments ----------
\newtheorem{theorem}{Theorem}
\newtheorem{definition}{Definition}
\newtheorem{lemma}{Lemma}
\newtheorem{proposition}{Proposition}
\newtheorem{remark}{Remark}
\newtheorem{example}{Example}
\newtheorem{conjecture}{Conjecture}

% ---------- Commands ----------
\newcommand{\Meas}{\mathbf{Meas}}
\newcommand{\Set}{\mathbf{Set}}
\newcommand{\Prob}{\mathbf{Prob}}
\newcommand{\Kl}{\mathrm{Kl}}
\newcommand{\Pcal}{\mathcal{P}}
\newcommand{\Xcal}{\mathcal{X}}
\newcommand{\Ycal}{\mathcal{Y}}
\newcommand{\Zcal}{\mathcal{Z}}
\newcommand{\Fcal}{\mathcal{F}}
\newcommand{\Bcal}{\mathcal{B}}
\newcommand{\Acal}{\mathcal{A}}
\newcommand{\Qcal}{\mathcal{Q}}
\newcommand{\Ocal}{\mathcal{O}}
\newcommand{\id}{\mathrm{id}}
\newcommand{\op}{\mathrm{op}}
\newcommand{\pr}{\mathrm{pr}}
\newcommand{\varlim}{\varprojlim}

\title{
The Observable Dynamics Program:\\
A Categorical Reading
}

\author{Patrick S. Mahon}

\date{Working note --- \today}

\begin{document}

\maketitle

\begin{abstract}
This note develops the Observable Dynamics Program in categorical language,
as a direct companion to the classical picture in \emph{categorical\_classical.tex}.
The classical construction is algebra-first: a probability measure is freely
chosen and marginals are derived from it.  The program is diagram-first: the
measure is derived from the coherence of a system of marginals.  We trace this
inversion through the same categorical layers --- $\Meas$, the Giry monad
$\Pcal$, its Kleisli category $\Kl(\Pcal)$ of Markov kernels, and the
extension theorem --- and identify where the program's arguments must go and
what remains open.  The inversion amounts to asking whether the forgetful
functor from probability spaces to compatible marginal systems has a left
adjoint, and under what conditions.
\end{abstract}

\tableofcontents

%%%%%%%%%%%%%%%%%%%%%%%%%%%%%%%%%%%%%%%%%%%%%%%%%%%%%%%%%%%%%%%%%%%%%
\section{The inversion: diagram-first probability}
%%%%%%%%%%%%%%%%%%%%%%%%%%%%%%%%%%%%%%%%%%%%%%%%%%%%%%%%%%%%%%%%%%%%%

The classical construction of probability (see \emph{categorical\_classical.tex},
\S6) has a single free parameter: the choice of measure $P \in \Pcal X$ at
Step~3.  Everything else --- pushforwards, marginals, kernels, projective
limits --- is forced by the categorical structure of $\Meas$ and $\Kl(\Pcal)$.

The Observable Dynamics Program asks: can Step~3 be replaced by a derivation?
Specifically: given a projective system of measurable spaces and a compatible
family of probability measures on its finite-dimensional factors, is there a
\emph{canonical} probability measure on the projective limit that is forced
by the coherence data, rather than freely chosen?

The program's answer is yes --- under conditions that we now identify
categorically.

\begin{center}
\begin{tabular}{ll}
\hline
\textbf{Classical (algebra-first)} & \textbf{Program (diagram-first)} \\
\hline
Choose $P \in \Pcal X$ freely & Derive $P$ from cocone data \\
Marginals follow from $P$ & $P$ follows from marginals \\
Monad algebra assumed & Monad algebra forced \\
Step 3 is a free parameter & Step 3 is a theorem \\
\hline
\end{tabular}
\end{center}

%%%%%%%%%%%%%%%%%%%%%%%%%%%%%%%%%%%%%%%%%%%%%%%%%%%%%%%%%%%%%%%%%%%%%
\section{Query systems as diagrams in $\Kl(\Pcal)$}
\label{sec:query-system}
%%%%%%%%%%%%%%%%%%%%%%%%%%%%%%%%%%%%%%%%%%%%%%%%%%%%%%%%%%%%%%%%%%%%%

\begin{definition}[Query system]\label{def:query-system}
A \emph{query system} $(\Qcal, \preceq)$ consists of:
\begin{itemize}
  \item A directed preorder $(I, \leq)$ of \emph{query indices}.
  \item For each $i \in I$, a measurable space $(O_i, \Ocal_i)$ --- the
    \emph{outcome space} of query $i$.
  \item For each $i \leq j$, a measurable function $\pi_{ij} : O_j \to O_i$
    --- the \emph{refinement map} from finer query $j$ to coarser query $i$
    --- satisfying $\pi_{ii} = \id$ and $\pi_{ij} \circ \pi_{jk} = \pi_{ik}$
    for $i \leq j \leq k$.
\end{itemize}
\end{definition}

\begin{remark}[Query systems as functors]
A query system is a functor $F : I^\op \to \Meas$, where $I$ is the preorder
category with objects the query indices and morphisms the order relations.
This is precisely a \emph{projective system} in $\Meas$.

The refinement maps $\pi_{ij}$ are deterministic, so they embed into
$\Kl(\Pcal)$ via the Dirac kernel $\hat\pi_{ij}(o, \cdot) = \delta_{\pi_{ij}(o)}$.
The query system therefore lives in $\Meas$, but it is a diagram in $\Kl(\Pcal)$
under this embedding.  We work in $\Kl(\Pcal)$ because the morphisms we
ultimately care about --- compatible marginals and extension kernels --- are
genuinely stochastic.
\end{remark}

\begin{definition}[Realization space]
The \emph{realization space} of a query system is the projective limit
\[
  \Omega = \varlim_{i \in I} O_i
  = \left\{ (o_i)_{i \in I} \in \prod_i O_i \;\middle|\;
    \pi_{ij}(o_j) = o_i \text{ for all } i \leq j \right\}
\]
equipped with the subspace $\sigma$-algebra from the product.  The projections
$\pr_i : \Omega \to O_i$ are the canonical measurable maps.
\end{definition}

\begin{remark}[What $\Omega$ represents]
A point $\omega \in \Omega$ is a \emph{coherent tuple} of query outcomes: a
choice of answer to every query simultaneously, consistent across all
refinement relations.  This is the program's analogue of a point in the sample
space $X$ --- except $\Omega$ is constructed from the queries, not assumed.

In the delay query system (Paper~4 / \emph{observational\_probability.tex}),
$\Omega \subseteq X^{\mathbb{Z}}$ is the set of bi-infinite sensor streams
that are coherent under all delay-and-sample operations.
\end{remark}

%%%%%%%%%%%%%%%%%%%%%%%%%%%%%%%%%%%%%%%%%%%%%%%%%%%%%%%%%%%%%%%%%%%%%
\section{Compatible marginals as cocones}
\label{sec:cocones}
%%%%%%%%%%%%%%%%%%%%%%%%%%%%%%%%%%%%%%%%%%%%%%%%%%%%%%%%%%%%%%%%%%%%%

\begin{definition}[Compatible marginals]\label{def:compat-marginals}
A \emph{compatible family of probability measures} for a query system
$(\{O_i\}, \{\pi_{ij}\})$ is a collection $\{\nu_i \in \Pcal O_i\}_{i \in I}$
satisfying
\[
  (\pi_{ij})_\sharp \nu_j = \nu_i \qquad \text{for all } i \leq j.
\]
\end{definition}

\begin{remark}[Cocone interpretation]
In $\Kl(\Pcal)$, the compatible family $\{\nu_i\}$ is a \emph{cocone} over
the diagram $F : I^\op \to \Kl(\Pcal)$.  Specifically, each $\nu_i$ is a
morphism $1 \to O_i$ in $\Kl(\Pcal)$ (a stochastic map from the terminal
object to $O_i$), and the compatibility condition says that these morphisms
commute with the refinement maps:
\[
\begin{tikzcd}
  1 \arrow[r, "\nu_j"] \arrow[dr, "\nu_i"'] & O_j \arrow[d, "\hat\pi_{ij}"] \\
  & O_i
\end{tikzcd}
\]
The apex of the cocone is (if it exists) a measure $P$ on $\Omega$ with
$(\pr_i)_\sharp P = \nu_i$ for all $i$.

The program's central question is: \emph{does this cocone have a canonical
apex?}
\end{remark}

\begin{remark}[Classical vs.\ program direction]
Classically, the apex $P$ is freely chosen first, and the cocone data
$\{\nu_i\}$ are derived as pushforwards $(\pr_i)_\sharp P$.  This is a
\emph{cone}: $P$ is the apex and the $\nu_i$ are its projections.

The program runs in the opposite direction: the cocone data $\{\nu_i\}$ are
given (from observations), and the apex $P$ is derived.  This is the
categorical form of the epistemological inversion: the global measure is
constituted by the coherence of its parts, not given as a prior that
determines the parts.
\end{remark}

%%%%%%%%%%%%%%%%%%%%%%%%%%%%%%%%%%%%%%%%%%%%%%%%%%%%%%%%%%%%%%%%%%%%%
\section{The extension theorem as an adjunction}
\label{sec:adjunction}
%%%%%%%%%%%%%%%%%%%%%%%%%%%%%%%%%%%%%%%%%%%%%%%%%%%%%%%%%%%%%%%%%%%%%

We now state the program's main theorem --- the Observational Extension
Theorem of Paper~1 --- in categorical language.

\begin{definition}[Categories for the adjunction]
Define two categories:
\begin{itemize}
  \item $\Prob(\Omega)$: the category whose objects are probability measures
    $P \in \Pcal\Omega$ and whose morphisms are measurable maps of probability
    spaces over $\Omega$.  (For route-finding purposes, we may take this to be
    a discrete category: objects are measures, no non-trivial morphisms.)
  \item $\mathrm{Compat}(\{O_i\})$: the category of compatible families
    $\{\nu_i\}$ over the query system, with morphisms being families of
    measure-preserving maps compatible with the refinements.
\end{itemize}
\end{definition}

\begin{definition}[The forgetful functor]
The \emph{marginalisation functor}
\[
  \Phi : \Prob(\Omega) \to \mathrm{Compat}(\{O_i\})
\]
sends a probability measure $P$ on $\Omega$ to the compatible family of its
marginals:
\[
  \Phi(P) = \left\{ (\pr_i)_\sharp P \right\}_{i \in I}.
\]
This functor is always well-defined: pushforward is functorial and the
compatibility condition holds automatically.
\end{definition}

\begin{theorem}[Extension theorem --- adjunction form]\label{thm:extension-adjunction}
Under the conditions of Paper~1 (query system sequentially upper-directed and
realizable), the functor $\Phi$ is an \emph{isomorphism of categories}:
every compatible family $\{\nu_i\}$ extends uniquely to a measure $P$ on
$\Omega$, and $\Phi(P) = \{\nu_i\}$.

In particular, $\Phi$ has a two-sided inverse (the extension functor), so
the adjunction $\mathrm{Ext} \dashv \Phi$ is in fact an equivalence.
\end{theorem}

\begin{proof}[Proof (sketch)]
The uniqueness direction --- that $\Phi$ is injective --- is the
Observational Determination Theorem (Paper~1, Theorem~5.2): two measures
on $\Omega$ with the same marginals on all query outcomes are equal, since
the query $\sigma$-algebra generates $\Bcal_\Omega$.

The existence direction --- that $\Phi$ has a left inverse --- is the
Observational Extension Theorem (Paper~1, Theorem~6.1): the
Carathéodory extension of the finitely additive content defined by the
$\{\nu_i\}$ on cylinder sets is $\sigma$-additive, hence extends to a
measure $P$ on $\Omega$.  The key steps are:
\begin{enumerate}
  \item Define $\mu_0$ on cylinder sets via $\mu_0(\pr_i^{-1}(A)) = \nu_i(A)$.
  \item Show $\mu_0$ is a well-defined finitely additive content
    (using compatibility and upper-directedness of the query system).
  \item Show $\mu_0$ is $\sigma$-subadditive (using realizability to pass
    from empty cylinder intersections to empty intersections in $\Omega$).
  \item Apply the Carathéodory extension theorem to obtain $P$.
\end{enumerate}
The Lean~4 formalization in \texttt{QuerySystem.lean} proves this in full
generality; see \texttt{observational\_extension} (line~1450ff).
\end{proof}

\begin{remark}[The Polish replacement]
The classical Kolmogorov extension theorem obtains existence by a topological
route: Polish spaces are regular enough that Prokhorov's tightness theorem
guarantees inner regularity, which in turn forces $\sigma$-additivity.

The program obtains existence by a purely measure-theoretic route: the
realizability condition replaces inner regularity.  Realizability says that
the projections $\pr_i : \Omega \to O_i$ are surjective, so that the
projective limit is ``full enough'' that vanishing marginals imply vanishing
cylinder content.

This is why the two approaches are complementary rather than redundant: the
program's route works in settings where the spaces need not be Polish (any
standard measurable space with the right query structure), and the
relationship to the Prokhorov/tightness route is precisely what the
Conjecture~2 / CFA work on the \texttt{equicontinuity-cfa-conjecture} branch
is probing.
\end{remark}

%%%%%%%%%%%%%%%%%%%%%%%%%%%%%%%%%%%%%%%%%%%%%%%%%%%%%%%%%%%%%%%%%%%%%
\section{The minimal predictive query in $\Kl(\Pcal)$}
\label{sec:predictive}
%%%%%%%%%%%%%%%%%%%%%%%%%%%%%%%%%%%%%%%%%%%%%%%%%%%%%%%%%%%%%%%%%%%%%

We now identify the objects of Paper~2 as morphisms in $\Kl(\Pcal)$.

\begin{definition}[Delay query system]
Fix a measurable sensor alphabet $(X, \Xcal)$ and work on the path space
$X^{\mathbb{Z}}$ with the product $\sigma$-algebra.  The \emph{delay query
system} has:
\begin{itemize}
  \item \emph{Queries}: pairs $(d, \tau) \in \mathbb{N}_{>0} \times \mathbb{N}_{>0}$
    with outcome space $O_{d,\tau} = X^d$ and evaluation map
    \[
      \mathrm{ev}_{d,\tau} : X^{\mathbb{Z}} \to X^d,
      \qquad
      \omega \mapsto (\omega_0, \omega_{-\tau}, \omega_{-2\tau}, \ldots, \omega_{-(d-1)\tau}).
    \]
  \item \emph{Refinement maps}: $\pi_{(d,\tau),(d',\tau')} : X^{d'} \to X^d$
    defined when $(d,\tau) \preceq (d',\tau')$, projecting to the
    corresponding coordinates.
\end{itemize}
The realization space is $\Omega \subseteq X^{\mathbb{Z}}$, the set of
streams coherent under all delay evaluations.
\end{definition}

\begin{definition}[Predictive kernel as Kleisli morphism]
Given a compatible family $\{\nu_{d,\tau}\}$ and the extended measure $P$ on
$\Omega$ (from Theorem~\ref{thm:extension-adjunction}), the \emph{predictive
kernel} at time $t$ is a morphism in $\Kl(\Pcal)$:
\[
  K_t : \Omega \to O_F, \qquad
  K_t(\omega, \cdot) = P\bigl((\sigma^t \omega \in \cdot) \mid \Fcal_{-\infty}^0\bigr),
\]
where $\sigma$ is the shift map and $O_F = X^{\mathbb{N}}$ is the future
observable space.  The \emph{minimal predictive query} $Q_*$ is the
$\sigma$-algebra generated by $\{K_t\}_{t \geq 0}$ --- the coarsest
measurable structure on $\Omega$ from which all future conditional
distributions can be computed.
\end{definition}

\begin{remark}[Semigroup in $\Kl(\Pcal)$]
The Markov property of $\{K_t\}$ is exactly closure of this collection under
Chapman--Kolmogorov composition in $\Kl(\Pcal)$:
\[
  K_{s+t} = K_t \circ K_s \qquad \text{(equality in } \Kl(\Pcal)\text{)}.
\]
The semigroup property of the predictive operator (Paper~2, Theorem~5.2) is
associativity of composition in $\Kl(\Pcal)$.  The operator-theoretic content
of Paper~3 (resolvent, spectral theory) is the analysis of this semigroup
acting on $L^2(\Omega, P)$.
\end{remark}

\begin{remark}[The Rose condition as isomorphism]
Two delay embeddings (via different $(d, \tau)$ parameters) are
\emph{observationally equivalent} if their induced predictive sub-diagrams
are isomorphic in $\Kl(\Pcal)$.  The Rose condition --- $D_f(\Gamma_Y \| \Gamma_Z) = 0$
--- is the information-theoretic criterion for this isomorphism.  It says that
no $f$-divergence can distinguish the two predictive distributions, which in
$\Kl(\Pcal)$ means the two kernels are equal as morphisms.
\end{remark}

%%%%%%%%%%%%%%%%%%%%%%%%%%%%%%%%%%%%%%%%%%%%%%%%%%%%%%%%%%%%%%%%%%%%%
\section{Open problems as categorical questions}
\label{sec:open}
%%%%%%%%%%%%%%%%%%%%%%%%%%%%%%%%%%%%%%%%%%%%%%%%%%%%%%%%%%%%%%%%%%%%%

The program's open problems now have precise categorical formulations.

\begin{conjecture}[Conjecture~2 --- $\sigma$-additivity from coherence alone]
\label{conj:sigma-additivity}
There exist purely algebraic conditions on the query system --- without
realizability or Polish hypotheses --- that force the finitely additive content
$\mu_0$ to be $\sigma$-additive.

\emph{Categorically}: what conditions on a diagram $F : I^\op \to \Kl(\Pcal)$
guarantee that the colimit exists in $\Pcal\Omega$?  This is a cocompleteness
question for $\Kl(\Pcal)$ restricted to diagrams of a particular shape.

\emph{Connection to equicontinuity}: the CFA (Compact Family Argument) work
on the \texttt{equicontinuity-cfa-conjecture} branch is developing the
Prokhorov-type tightness argument for this.  If successful, it would give a
topological route to cocompleteness analogous to the classical case, and
potentially feed into a Mathlib formalization of the Kolmogorov extension
theorem.
\end{conjecture}

\begin{conjecture}[Stopping Point~1 --- measurability from discriminability]
\label{conj:measurability}
The $\sigma$-algebra $\Xcal$ on the sensor alphabet $X$ should be derivable
from primitive discriminative structure: the ability of the observer to
distinguish finite families of events.

\emph{Categorically}: is there a category more primitive than $\Meas$ ---
perhaps of \emph{locales}, \emph{domains}, or \emph{topological spaces} ---
whose image in $\Meas$ under a forgetful functor gives the right
$\sigma$-algebras?  The $\sigma$-algebra would then be the \emph{sheaf of
reportable distinctions} on a suitable site, and the assumption of $(X, \Xcal)$
in the current program would be replaced by a construction.

\emph{Connection to Boole}: Boolean algebras of propositions (Boole's algebra
of thought) are the finite discriminative structure; the $\sigma$-algebra is
their countable closure.  The conjecture is that any system capable of
coherent countable probabilistic reasoning must already have this closure.
\end{conjecture}

\begin{remark}[The adjunction gap]
Theorem~\ref{thm:extension-adjunction} states the extension theorem as an
isomorphism of categories.  The full adjunction statement --- identifying the
unit and counit, showing they satisfy the triangle identities, and characterising
the fixpoints --- has not been written out precisely.  This is a concrete
mathematical task: define the two categories properly, write down the functors
$\Phi$ and $\mathrm{Ext}$, and verify the adjunction axioms.  It would give
the cleanest categorical statement of what Paper~1 proves.
\end{remark}

%%%%%%%%%%%%%%%%%%%%%%%%%%%%%%%%%%%%%%%%%%%%%%%%%%%%%%%%%%%%%%%%%%%%%
\section{Comparison table: classical vs.\ program}
\label{sec:comparison}
%%%%%%%%%%%%%%%%%%%%%%%%%%%%%%%%%%%%%%%%%%%%%%%%%%%%%%%%%%%%%%%%%%%%%

\begin{center}
\renewcommand{\arraystretch}{1.3}
\begin{tabular}{lll}
\hline
\textbf{Step} & \textbf{Classical} & \textbf{Program} \\
\hline
1. Starting point & Sample space $(X, \Xcal)$ assumed & Query system $(\{O_i\}, \{\pi_{ij}\})$ given \\
2. Ambient category & $\Meas$ and $\Kl(\Pcal)$ & $\Meas$ and $\Kl(\Pcal)$ (same) \\
3. Measure & Free choice $P : 1 \to \Pcal X$ & Derived: $P = \mathrm{Ext}(\{\nu_i\})$ \\
4. Marginals & Derived from $P$ & Given as input (observations) \\
5. Limit space & $X$ is assumed & $\Omega = \varlim O_i$ is constructed \\
6. Extension & Kolmogorov (Polish) & Paper~1 (sequential, realizable) \\
7. Dynamics & Markov kernel $T : X \to X$ in $\Kl(\Pcal)$ & Predictive kernel $K_t : \Omega \to O_F$ \\
8. Semigroup & Chapman--Kolmogorov & Composition in $\Kl(\Pcal)$ (same) \\
9. Objectivity & Ground truth assumed & Limit of coherent query perspectives \\
\hline
\end{tabular}
\end{center}

Row~3 is the inversion.  Rows~2 and~8 show that the categorical machinery
is the same in both directions --- the difference is not in the tools but in
the direction of travel.

Row~9 deserves emphasis: the program does not eliminate objectivity but
\emph{derives} it.  When all queries cohere --- when the compatible family
$\{\nu_i\}$ is globally consistent --- the extension $P$ is the unique
probability measure on $\Omega$ that everyone with access to all queries
would agree on.  This is the mathematical form of the phenomenological claim
that objectivity is constituted by the coherence of perspectives, not
presupposed beneath them.

%%%%%%%%%%%%%%%%%%%%%%%%%%%%%%%%%%%%%%%%%%%%%%%%%%%%%%%%%%%%%%%%%%%%%
\section{Mathlib availability for the program's picture}
\label{sec:mathlib}
%%%%%%%%%%%%%%%%%%%%%%%%%%%%%%%%%%%%%%%%%%%%%%%%%%%%%%%%%%%%%%%%%%%%%

\begin{itemize}
  \item \textbf{Available}: \texttt{CategoryTheory.Monad.Kleisli} gives the
    general Kleisli construction.  Instantiated at the Giry monad, morphisms
    are Markov kernels (\texttt{Probability.Kernel}), which are extensively
    developed.

  \item \textbf{Available}: \texttt{CategoryTheory.Monad.Algebra} gives
    the Eilenberg--Moore category and the free/forgetful adjunction.  The
    query system adjunction (Theorem~\ref{thm:extension-adjunction}) is a
    specific instance of this pattern.

  \item \textbf{Available}: \texttt{MeasureTheory.Measure.GiryMonad} for
    the monad laws; \texttt{MeasureTheory.Constructions.Pi} for finite
    products; \texttt{MeasureTheory.Constructions.Projective} for
    projective families (partial).

  \item \textbf{Built in-house}: \texttt{QuerySystem.lean} provides
    \texttt{observational\_extension} --- the main theorem --- as new Lean~4
    infrastructure not yet in Mathlib.  This fills the Kolmogorov extension
    gap for the query system setting.

  \item \textbf{Gap}: the precise adjunction of
    Theorem~\ref{thm:extension-adjunction} (defining $\Prob(\Omega)$ and
    $\mathrm{Compat}(\{O_i\})$ as Lean~4 categories and verifying the
    adjunction axioms) has not been formalized.  This is the next categorical
    formalization target.
\end{itemize}

\end{document}
